% ============================================================
% 04_ADA.TEX
% Sistema Adaptativo de Classificação Multicritério
% de Fluidos de Corte
%
% FINALIDADE
%
% Apresentar a Análise Descritiva/Exploratória dos Dados:
%
% - comportamento individual das variáveis;
% - comparação entre fluidos;
% - comparação entre condições;
% - associações entre variáveis;
% - estrutura multivariada;
% - identificação de redundâncias;
% - preparação para a seleção dos critérios.
%
% Este arquivo constitui apenas a BASE da apresentação.
%
% Estatísticas, figuras e conclusões deverão ser inseridas
% pelos integrantes somente após a execução das análises em R.
%
% ============================================================


% ============================================================
% FRAME 1 — OBJETIVO DA ADA
% ============================================================

\begin{frame}{Análise exploratória dos dados}

    \begin{block}{Objetivo}
        Compreender a estrutura dos dados antes da definição da
        matriz multicritério.
    \end{block}

    \vspace{0.3cm}

    A análise exploratória deverá investigar:

    \begin{itemize}
        \item distribuição das variáveis;
        \item diferenças entre fluidos;
        \item comportamento entre condições experimentais;
        \item variabilidade;
        \item valores extremos;
        \item associações entre variáveis;
        \item redundância potencial;
        \item estrutura multivariada.
    \end{itemize}

    \vspace{0.2cm}

    \begin{alertblock}{Importante}
        A ADA serve para compreender os dados e fundamentar decisões
        posteriores. Ela não produz, isoladamente, o ranking final.
    \end{alertblock}

\end{frame}


% ============================================================
% FRAME 2 — ANÁLISE UNIVARIADA
% ============================================================

\begin{frame}{Análise univariada}

    Cada variável será inicialmente examinada de forma individual.

    \vspace{0.3cm}

    Quando adequadas à natureza da variável, poderão ser utilizadas:

    \begin{itemize}
        \item média;
        \item mediana;
        \item mínimo e máximo;
        \item quartis;
        \item desvio-padrão;
        \item amplitude;
        \item frequência de valores ausentes;
        \item distribuição gráfica.
    \end{itemize}

    \vspace{0.3cm}

    \begin{block}{Objetivo}
        Identificar escala, dispersão, assimetria, baixa variabilidade
        e valores que mereçam investigação adicional.
    \end{block}

\end{frame}


% ============================================================
% FRAME 3 — COMPARAÇÃO ENTRE FLUIDOS
% ============================================================

\begin{frame}{Comparação entre fluidos}

    A ADA deverá comparar o comportamento das alternativas em cada
    característica relevante.

    \vspace{0.3cm}

    Poderão ser utilizados, conforme os dados:

    \begin{itemize}
        \item boxplots;
        \item dotplots;
        \item gráficos de pontos;
        \item gráficos por condição;
        \item tabelas descritivas compactas.
    \end{itemize}

    \vspace{0.3cm}

    \begin{block}{Pergunta exploratória}
        Quais características diferenciam efetivamente os fluidos?
    \end{block}

    \vspace{0.2cm}

    \begin{alertblock}{Cuidado}
        Uma alternativa favorável em uma variável não deve ser
        classificada como a melhor alternativa global.
    \end{alertblock}

\end{frame}


% ============================================================
% FRAME 4 — CONDIÇÕES EXPERIMENTAIS
% ============================================================

\begin{frame}{Comportamento entre condições}

    Quando uma variável estiver disponível em diferentes condições,
    será necessário avaliar sua estabilidade.

    \vspace{0.3cm}

    A análise deverá investigar:

    \begin{itemize}
        \item mudança de magnitude;
        \item mudança de posição relativa;
        \item estabilidade;
        \item possíveis interações;
        \item inversão de desempenho;
        \item necessidade ou não de agregação.
    \end{itemize}

    \vspace{0.3cm}

    \begin{block}{Implicação}
        Um valor médio pode ocultar comportamentos distintos entre
        condições experimentais.
    \end{block}

\end{frame}


% ============================================================
% FRAME 5 — VALORES EXTREMOS
% ============================================================

\begin{frame}{Valores extremos}

    Valores extremos deverão ser tratados inicialmente como:

    \begin{center}
        \textbf{informação a investigar}
    \end{center}

    e não como:

    \begin{center}
        \textbf{erro a remover automaticamente}.
    \end{center}

    \vspace{0.3cm}

    Para cada caso relevante será necessário avaliar:

    \begin{itemize}
        \item origem do valor;
        \item consistência com a fonte;
        \item plausibilidade técnica;
        \item influência nas estatísticas;
        \item impacto sobre as etapas posteriores.
    \end{itemize}

\end{frame}


% ============================================================
% FRAME 6 — ASSOCIAÇÃO ENTRE VARIÁVEIS
% ============================================================

\begin{frame}{Associação entre variáveis}

    A ADA bivariada deverá investigar relações entre critérios
    candidatos.

    \vspace{0.3cm}

    Poderão ser utilizadas:

    \begin{itemize}
        \item matriz de dispersão;
        \item correlação de Pearson;
        \item correlação de Spearman, quando pertinente;
        \item visualização por pares;
        \item análise técnica das variáveis.
    \end{itemize}

    \vspace{0.3cm}

    \begin{block}{Objetivos}
        Identificar possíveis redundâncias e trade-offs antes da
        definição da matriz de decisão.
    \end{block}

    \vspace{0.2cm}

    \begin{alertblock}{Interpretação}
        Correlação indica associação e não implica causalidade.
    \end{alertblock}

\end{frame}


% ============================================================
% FRAME 7 — REDUNDÂNCIA
% ============================================================

\begin{frame}{Redundância entre critérios}

    Duas variáveis muito semelhantes podem representar praticamente
    a mesma informação.

    \vspace{0.3cm}

    \begin{block}{Risco}
        Utilizar simultaneamente critérios redundantes pode provocar
        super-representação de uma mesma dimensão no sistema
        multicritério.
    \end{block}

    \vspace{0.3cm}

    A investigação deverá considerar:

    \begin{itemize}
        \item intensidade da associação;
        \item significado técnico;
        \item forma de cálculo;
        \item dependência matemática;
        \item origem das variáveis.
    \end{itemize}

    \vspace{0.2cm}

    \centering
    \textit{Correlação elevada, isoladamente, não determina exclusão.}

\end{frame}


% ============================================================
% FRAME 8 — CORRELAÇÃO POSITIVA E NEGATIVA
% ============================================================

\begin{frame}{Interpretação das correlações}

    \begin{columns}[T]

        \begin{column}{0.48\textwidth}

            \begin{block}{Correlação positiva}
                Pode indicar que duas características variam em
                direções semelhantes e potencialmente compartilham
                informação.
            \end{block}

        \end{column}


        \begin{column}{0.48\textwidth}

            \begin{block}{Correlação negativa}
                Pode representar conflito ou trade-off entre
                características e não necessariamente redundância.
            \end{block}

        \end{column}

    \end{columns}

    \vspace{0.5cm}

    Essa distinção será particularmente importante posteriormente no
    método CRITIC, que utiliza as relações entre os critérios para
    quantificar conteúdo informacional.

\end{frame}


% ============================================================
% FRAME 9 — ANÁLISE MULTIVARIADA
% ============================================================

\begin{frame}{Estrutura multivariada}

    Quando pertinente, técnicas multivariadas poderão auxiliar na
    compreensão da estrutura conjunta dos dados.

    \vspace{0.3cm}

    Exemplos possíveis:

    \begin{itemize}
        \item análise de componentes principais;
        \item mapas de calor;
        \item análise de agrupamentos;
        \item representações em espaço reduzido.
    \end{itemize}

    \vspace{0.3cm}

    \begin{block}{Finalidade}
        Investigar padrões conjuntos, proximidade entre alternativas e
        relações entre variáveis.
    \end{block}

    \vspace{0.2cm}

    \begin{alertblock}{Limite}
        Essas técnicas terão caráter exploratório e não substituirão
        o método multicritério utilizado para o ranking.
    \end{alertblock}

\end{frame}


% ============================================================
% FRAME 10 — PCA
% ============================================================

\begin{frame}{Análise de Componentes Principais}

    Caso a PCA seja utilizada, ela poderá ajudar a visualizar:

    \begin{itemize}
        \item estrutura de correlação;
        \item grupos de variáveis;
        \item proximidade entre fluidos;
        \item variáveis associadas às principais direções de variação.
    \end{itemize}

    \vspace{0.3cm}

    \begin{block}{A preencher após execução}
        \begin{itemize}
            \item Variância explicada pela PC1:
                  \texttt{TO\_BE\_FILLED}
            \item Variância explicada pela PC2:
                  \texttt{TO\_BE\_FILLED}
            \item Principais padrões:
                  \texttt{TO\_BE\_FILLED}
        \end{itemize}
    \end{block}

    \vspace{0.2cm}

    \centering
    \textit{Somente incluir este slide se a PCA for efetivamente utilizada.}

\end{frame}


% ============================================================
% FRAME 11 — PCA NÃO É RANKING
% ============================================================

\begin{frame}{PCA não define o melhor fluido}

    \begin{alertblock}{Cuidado de interpretação}
        A posição de uma alternativa em uma componente principal não
        representa automaticamente sua qualidade global.
    \end{alertblock}

    \vspace{0.4cm}

    A PCA busca representar:

    \begin{center}
        \textbf{variabilidade dos dados}
    \end{center}

    e não:

    \begin{center}
        \textbf{preferência multicritério}.
    \end{center}

    \vspace{0.3cm}

    Portanto:

    \[
    \text{PCA}
    \neq
    \text{ranking TOPSIS}
    \]

\end{frame}


% ============================================================
% FRAME 12 — VARIÁVEIS CONSTANTES
% ============================================================

\begin{frame}{Capacidade discriminatória}

    Uma característica utilizada como critério precisa contribuir para
    distinguir as alternativas.

    \vspace{0.3cm}

    Variáveis poderão ser classificadas como:

    \begin{itemize}
        \item discriminatórias;
        \item pouco variáveis;
        \item constantes;
        \item diagnósticas;
        \item redundantes;
        \item inadequadas como critério.
    \end{itemize}

    \vspace{0.3cm}

    Para uma variável constante:

    \[
    \operatorname{Var}(X)=0
    \]

    ela não possui capacidade de discriminar os fluidos naquele conjunto
    analisado.

\end{frame}


% ============================================================
% FRAME 13 — DA ADA AOS CRITÉRIOS
% ============================================================

\begin{frame}{Da exploração à definição dos critérios}

    \begin{block}{Fluxo}
        \centering

        Variáveis disponíveis
        $\rightarrow$
        ADA
        $\rightarrow$
        Redundância
        $\rightarrow$
        Conformidade
        $\rightarrow$
        Critérios candidatos
        $\rightarrow$
        Critérios aprovados
    \end{block}

    \vspace{0.4cm}

    A definição final dos critérios deverá considerar conjuntamente:

    \begin{itemize}
        \item informação estatística;
        \item significado técnico;
        \item disponibilidade dos dados;
        \item variabilidade;
        \item redundância;
        \item direção de preferência;
        \item possibilidade de interpretação.
    \end{itemize}

\end{frame}


% ============================================================
% FRAME 14 — RESULTADOS FUTUROS DA ADA
% ============================================================

\begin{frame}{Principais achados exploratórios}

    \begin{block}{Achado 1}
        \texttt{TO\_BE\_FILLED}
    \end{block}

    \begin{block}{Achado 2}
        \texttt{TO\_BE\_FILLED}
    \end{block}

    \begin{block}{Achado 3}
        \texttt{TO\_BE\_FILLED}
    \end{block}

    \vspace{0.3cm}

    \centering

    \textit{Preencher somente após a execução e interpretação das análises
    pelos integrantes.}

\end{frame}


% ============================================================
% FRAME 15 — RESULTADOS QUE DEVEM SER LEVADOS ADIANTE
% ============================================================

\begin{frame}{Síntese para a etapa multicritério}

    Após a ADA, deverão estar definidos ou documentados:

    \begin{itemize}
        \item variáveis com capacidade discriminatória;
        \item associações importantes;
        \item redundâncias potenciais;
        \item variáveis constantes ou quase constantes;
        \item comportamento entre condições;
        \item valores extremos relevantes;
        \item limitações;
        \item critérios candidatos.
    \end{itemize}

    \vspace{0.3cm}

    \begin{block}{Próxima etapa}
        Transformar essas evidências em uma matriz de decisão
        tecnicamente justificável.
    \end{block}

\end{frame}


% ============================================================
% FIGURAS FUTURAS
% ============================================================
%
% Os integrantes poderão substituir slides mais textuais por
% figuras produzidas em R, por exemplo:
%
% - boxplot de variável relevante;
% - gráfico por condição;
% - matriz de correlação;
% - heatmap;
% - biplot da PCA;
% - gráfico de dispersão;
%
% Priorizar poucas figuras com alto valor explicativo.
%
% Não é necessário produzir nenhuma agora.
%
% ============================================================


% ============================================================
% MATERIAL EXTERNO
% ============================================================
%
% Caso a professora forneça:
%
% - interpretação oficial de variável;
% - significado técnico;
% - referência;
% - critério obrigatório;
%
% esse material deverá ser incorporado antes de tomar decisões
% sobre exclusão ou uso de critérios.
%
% ============================================================


% ============================================================
% ESTADO ATUAL
% ============================================================
%
% ADA_SECTION_STATUS = BASE_PREPARED
%
% UNIVARIATE_ADA_STATUS = NOT_EXECUTED
% BIVARIATE_ADA_STATUS = NOT_EXECUTED
% MULTIVARIATE_ADA_STATUS = NOT_EXECUTED
% PCA_STATUS = TO_BE_DECIDED
%
% ADA_MAIN_FINDINGS = TO_BE_FILLED
% POTENTIAL_REDUNDANCIES = TO_BE_FILLED
% CANDIDATE_CRITERIA = TO_BE_FILLED
%
% ============================================================